%!TEX root = ../template.tex
%%%%%%%%%%%%%%%%%%%%%%%%%%%%%%%%%%%%%%%%%%%%%%%%%%%%%%%%%%%%%%%%%%%
%% chapter1.tex
%% Minimal chapter template — content removed, structure preserved
%%%%%%%%%%%%%%%%%%%%%%%%%%%%%%%%%%%%%%%%%%%%%%%%%%%%%%%%%%%%%%%%%%%

	ypeout{NT FILE chapter1.tex}%

% --- Chapter header (edit title and label) ---------------------------------
\chapter{Chapter Title}
\label{cha:chapter1}

% Graphics path kept (edit if you add chapter-specific figures)
\prependtographicspath{{Chapters/Figures/Covers/}}

% Epigraph configuration kept (remove or edit if not needed)
\epigraphfontsize{\small\itshape}
\setlength\epigraphwidth{12.5cm}
\setlength\epigraphrule{0pt}

% ------------------ Writing placeholders (replace these with your text) ----

\section{Motivation}
\label{sec:motivation}
% Write the motivation for this chapter here.

\section{Goals}
\label{sec:goals}
% List the goals or research questions for this chapter here.

\section{Report Structure}
\label{sec:report_structure}
% Give a short description of how this chapter fits into the report.

% ----- Helpful placeholders and notes --------------------------------------
% Figures: put chapter figures into Chapters/Figures/<this-chapter>/ and use
% \begin{figure} ... \end{figure}
%
% Tables: use standard LaTeX tabular or longtable for large tables.
%
% Citations: add \cite{key} entries to Bibliography/bibliography.bib as needed.
%
% Acronyms and glossary: add new acronyms in Chapters/acronyms.tex and use
% \gls{} and \glspl{} commands here.



